% ============================================================================
% RELATED: ALGEBRAIC STRUCTURE AND GROUP-EQUIVARIANT COMPUTATION
% ============================================================================
\subsection{Algebraic Structure and Group-Equivariant Computation}
\label{sec:related_algebraic_topology}

The Geometric Deep Learning programme formalized by Bronstein et al.\
\citep{bronstein2021geometric} demonstrates that many successful neural
architectures---CNNs, GNNs, Transformers---can be derived as special
cases of equivariant message-passing on domains with specific symmetry
groups.  Cohen and Welling's work on Group Equivariant Convolutional
Networks \citep{cohen2016group} provides the formal justification: when
the underlying data possesses a known symmetry, restricting
transformations to that symmetry group improves sample efficiency and
generalization.

The Value architecture relates to this programme through two algebraic
lanes.  The low-nibble Boolean instruction set forms a closed algebraic
structure: composing any two operations yields another operation
expressible as a sequence of truth-table applications.  When the same
instruction is applied uniformly across all bit positions of a frame, the
operation is translation-equivariant in the bit-position coordinate---
shifting the input by $k$ positions and then applying the operation
produces the same result as applying the operation and then shifting.
The high-nibble PGA lane adds a continuous group-action channel for
rigid transitions and reversals.

This equivariance is less expressive than learned group-equivariant
networks, which can adapt their feature maps to continuous symmetries of
the input domain.  The Value architecture instead provides fixed
algebraic structures---Boolean algebra over $\{0,1\}^{8192}$ and a
small PGA multivector algebra over selected 64-byte lanes---that are
deterministic and algebraically closed by construction.  The trade-off
is rigidity: the system does not learn arbitrary new symmetries from
data, but the symmetries it does possess are guaranteed rather than
approximate, and they can be exploited at the hardware level.

\paragraph{Non-Commutativity and Sequence Sensitivity.}
Although individual boolean operations are bitwise-commutative (AND is
commutative), sequences of operations applied to a frame are generally
\emph{non-commutative}: applying operation $A$ followed by operation $B$
to a frame typically produces a different result than applying $B$
followed by $A$.  This property allows the order in which frames
interact during Region mixing to encode structural information.  Two
sequences of the same frames, mixed in different orders, may yield
different final states.  In this sense, the non-commutativity of
composed Boolean and geometric operations serves a role analogous to
positional encoding in sequence models, without requiring a learned
position representation.
